\sheettitle
  {Preparatory notes}
  {From ambient $L^p$ spaces to probability representations}

\section*{0. The ambient measure is part of the object}

Let $(E,\mathcal A,\mu)$ be a measure space. For $1\le p<\infty$,
\[
L^p(\mu)
=
\left\{
f:E\to\mathbb R\text{ measurable}:
\int_E|f|^p\,d\mu<\infty
\right\}\big/\!\sim_\mu,
\]
where
\[
f\sim_\mu g
\quad\Longleftrightarrow\quad
f=g\quad\mu\text{-almost everywhere}.
\]
The norm is
\[
\|f\|_p
=
\left(\int_E|f|^p\,d\mu\right)^{1/p}.
\]
The space $L^\infty(\mu)$ consists of the equivalence classes of essentially bounded measurable functions, with
\[
\|f\|_\infty
=
\inf\{M\ge0:|f|\le M\ \mu\text{-almost everywhere}\}.
\]

Changing the ambient measure may change integrability, the norm, the null sets, and therefore the equivalence class represented by a fixed pointwise formula. On a probability space,
\[
X\in L^p(\mathbb P)
\quad\Longleftrightarrow\quad
\mathbb E[|X|^p]<\infty.
\]

\section*{1. Finite-measure embeddings}

\textbf{Proposition 1.}
Assume $\mu(E)<\infty$ and let $1\le p<q\le\infty$. Then
\[
L^q(\mu)\subseteq L^p(\mu),
\qquad
\|f\|_p
\le
\mu(E)^{1/p-1/q}\|f\|_q,
\]
with the convention $1/\infty=0$.

\textit{Proof.}
Suppose first that $q<\infty$. Hölder's inequality with conjugate exponents
\[
r=\frac qp,
\qquad
r'=\frac q{q-p}
\]
gives
\[
\begin{aligned}
\|f\|_p^p
&=
\int_E |f|^p\mathbf 1_E\,d\mu\\
&\le
\left(\int_E|f|^q\,d\mu\right)^{p/q}
\left(\int_E\mathbf 1_E\,d\mu\right)^{1-p/q}\\
&=
\|f\|_q^p\mu(E)^{1-p/q}.
\end{aligned}
\]
Taking the $p$-th root yields the estimate. If $q=\infty$, then
\[
\|f\|_p^p
\le
\|f\|_\infty^p\mu(E).
\]
This is the same estimate at the endpoint.

The hypothesis $\mu(E)<\infty$ enters through the integrability of $\mathbf 1_E$. No inclusion of this form is asserted on an arbitrary infinite measure space.

\textbf{Probability-space consequence.}
Since $\mathbb P(\Omega)=1$,
\[
\|X\|_p\le\|X\|_q
\qquad(1\le p<q\le\infty),
\]
and in particular
\[
L^\infty(\mathbb P)
\subseteq
L^2(\mathbb P)
\subseteq
L^1(\mathbb P).
\]
The three spaces serve different purposes: $L^1$ gives finite expectations and finite signed RN numerators; $L^\infty$ supplies indicators and bounded tests; $L^2$ legitimizes second moments and cross-products.

\section*{2. Two integrable products}

\textbf{Proposition 2.}
If $Y\in L^1(\mu)$ and $Z\in L^\infty(\mu)$, then
\[
YZ\in L^1(\mu),
\qquad
\|YZ\|_1
\le
\|Z\|_\infty\|Y\|_1.
\]

\textit{Proof.}
The estimate follows by integrating
\[
|YZ|
\le
\|Z\|_\infty |Y|
\quad\mu\text{-almost everywhere}.
\]

\textbf{Proposition 3.}
If $U,V\in L^2(\mu)$, then
\[
UV\in L^1(\mu),
\qquad
\|UV\|_1\le\|U\|_2\|V\|_2.
\]

\textit{Proof.}
This is the Cauchy--Schwarz inequality applied to $|U|$ and $|V|$.

The first proposition legitimizes bounded test functions against $L^1$ variables. The second legitimizes covariance and the cross-term in variance decompositions. Neither statement requires the Hilbert-space projection theorem.

\section*{3. From event identities to bounded tests}

\textbf{Proposition 4.}
Let $\mathcal G\subseteq\mathcal A$ be a sub-$\sigma$-field and let $f,g\in L^1(\mu)$. If
\[
\int_A f\,d\mu
=
\int_A g\,d\mu
\qquad\text{for every }A\in\mathcal G,
\]
then
\[
\int_E Zf\,d\mu
=
\int_E Zg\,d\mu
\]
for every bounded $\mathcal G$-measurable $Z$.

\textit{Proof.}
Linearity gives the identity first for bounded $\mathcal G$-measurable simple functions. Let $Z$ be bounded and $\mathcal G$-measurable. Choose simple functions $Z_n$ such that $Z_n\to Z$ pointwise and $|Z_n|\le\|Z\|_\infty$. Then
\[
|Z_n(f-g)|
\le
\|Z\|_\infty|f-g|\in L^1(\mu).
\]
Dominated convergence passes the identity to $Z$.

If $f$ and $g$ are themselves $\mathcal G$-measurable, taking $A=\{f>g\}$ and $A=\{g>f\}$ also shows
\[
f=g\quad\mu\text{-almost everywhere}.
\]

\section*{4. Density-generated measures}

\textbf{Proposition 5.}
If $h:E\to[0,\infty]$ is measurable, then
\[
(h\mu)(A)=\int_Ah\,d\mu
\]
defines a positive measure and $h\mu\ll\mu$. No integrability or $\sigma$-finiteness hypothesis is required for this forward construction.

\textit{Proof.}
For pairwise disjoint measurable sets $A_n$, monotone convergence applied to
\[
h\mathbf 1_{\cup_n A_n}
=
\sum_{n\ge1}h\mathbf 1_{A_n}
\]
gives countable additivity. If $\mu(A)=0$, then $\int_Ah\,d\mu=0$.

\textbf{Proposition 6.}
If $f\in L^1(\mu)$, then
\[
\nu_f(A)=\int_Af\,d\mu
\]
defines a finite signed measure, $\nu_f\ll\mu$, and
\[
|\nu_f|=|f|\mu.
\]
If $g\in L^1(\mu)$ and $\nu_f=\nu_g$, then
\[
f=g\quad\mu\text{-almost everywhere}.
\]

\textit{Proof.}
Write $f=f^+-f^-$. The positive measures $f^+\mu$ and $f^-\mu$ are finite and mutually singular, so they form the Jordan decomposition of $\nu_f$; hence $|\nu_f|=(f^++f^-)\mu=|f|\mu$. Absolute continuity is immediate.

For uniqueness, put $u=f-g$. The identity $\nu_f=\nu_g$ gives
\[
\int_Au\,d\mu=0
\qquad(A\in\mathcal A).
\]
Taking $A=\{u>0\}$ and $A=\{u<0\}$ shows that $u^+=u^-=0$ almost everywhere.

\section*{5. Change of measure}

\textbf{Proposition 7.}
Let $\nu=h\mu$ with $h\ge0$ measurable. For every nonnegative measurable $\varphi$,
\[
\int_E\varphi\,d\nu
=
\int_E\varphi h\,d\mu,
\]
where either side may equal $+\infty$. For signed $\varphi$, the same identity holds whenever
\[
\int_E|\varphi|\,d\nu
=
\int_E|\varphi|h\,d\mu
<\infty.
\]
Consequently,
\[
\varphi\in L^1(\nu)
\quad\Longleftrightarrow\quad
|\varphi|h\in L^1(\mu).
\]

\textit{Proof.}
The identity follows for indicators from the definition of $\nu$, then for nonnegative simple functions by linearity, and finally for arbitrary nonnegative functions by monotone convergence. Apply the result to positive and negative parts in the integrable signed regime.

Since $\nu\ll\mu$, every $\mu$-almost-everywhere identity is also a $\nu$-almost-everywhere identity. Absolute continuity alone does not imply that every element of $L^1(\mu)$ belongs to $L^1(\nu)$.

\section*{6. RN theorem boundary and signed extension}

\textbf{Radon--Nikodym theorem, positive standard regime.}
Let $\mu$ and $\rho$ be $\sigma$-finite positive measures on $(E,\mathcal A)$. If $\rho\ll\mu$, then there exists a measurable $h\ge0$ such that
\[
\rho=h\mu.
\]
The density $h=d\rho/d\mu$ is unique $\mu$-almost everywhere.

If $\mu$ is $\sigma$-finite and $\nu$ is a finite signed measure satisfying $\nu\ll\mu$, apply the positive theorem to the Jordan parts $\nu^+$ and $\nu^-$. Their densities $h^+$ and $h^-$ give
\[
\nu=(h^+-h^-)\mu.
\]
This is the legitimate signed extension; a positive-measure theorem is not silently applied to a signed numerator.

\section*{7. Four later uses}

\begin{enumerate}
  \item \textbf{General density, 14 July:} numerator $f\mu$, reference $\mu$, uniqueness $\mu$-almost everywhere.
  \item \textbf{Event conditioning, 19 July:} numerator $\mathbb P_B$, reference $\mathbb P$, explicit density $\mathbf 1_B/\mathbb P(B)$ when $\mathbb P(B)>0$.
  \item \textbf{State-space density, 24 July:} numerator $\mu_X=X_\#\mathbb P$, reference chosen on the state space, typically Lebesgue measure. A law is not a density before this choice and an absolute-continuity relation.
  \item \textbf{Conditional expectation, 26 July:} for $Y\in L^1(\mathbb P)$ and a sub-$\sigma$-field $\mathcal G\subseteq\mathcal F$, numerator $A\mapsto\mathbb E[Y\mathbf 1_A]$ on $\mathcal G$, reference $\mathbb P|_{\mathcal G}$, uniqueness $\mathbb P$-almost surely.
\end{enumerate}

The RN theorem consumes a numerator measure, a reference measure, absolute continuity, and the correct measure-theoretic regime. It is not a formal consequence of an $L^p$ embedding.
